\documentclass[letterpaper]{article} 
\usepackage[utf8]{inputenc}
\linespread{0.85}
\usepackage[T1]{fontenc}
\usepackage{amsmath}
\usepackage{amsfonts}
\usepackage{amssymb}
\usepackage{array}
\usepackage{booktabs}
\usepackage{hyperref}
\usepackage[version=4]{mhchem}
\usepackage{stmaryrd}
\usepackage{tikz}
\usepackage{amsmath}
\usepackage{graphicx}
\usepackage{capt-of}
\usepackage{lipsum}
\usepackage{fancyvrb}
\usepackage{tabularx}
\usepackage{listings}
\usepackage[export]{adjustbox}
\graphicspath{ {./images/} }
\usepackage[utf8]{inputenc}
\usepackage[english]{babel}
\usepackage{float}
\usepackage{lipsum}
\usepackage{graphicx}
\usepackage{float}
\usepackage[margin=0.7in]{geometry}
\usepackage{amsmath}
\usepackage{graphicx}
\usepackage{capt-of}
\usepackage{tcolorbox}
\usepackage{lipsum}
\usepackage{graphicx}
\usepackage{float}
\usepackage{listings}
\usepackage{hyperref} 
\usepackage{xcolor} % For custom colors
\lstset{
	language=Python,                % Choose the language (e.g., Python, C, R)
	basicstyle=\ttfamily\small, % Font size and type
	keywordstyle=\color{blue},  % Keywords color
	commentstyle=\color{gray},  % Comments color
	stringstyle=\color{red},    % String color
	numbers=left,               % Line numbers
	numberstyle=\tiny\color{gray}, % Line number style
	stepnumber=1,               % Numbering step
	breaklines=true,            % Auto line break
	backgroundcolor=\color{black!5}, % Light gray background
	frame=single,               % Frame around the code
}
\usepackage{float}
\usepackage[]{amsthm} % use \begin{proof}
\usepackage[]{amssymb} %gives us the character \varnothing
\usepackage{algorithm}
\usepackage{algpseudocode}

	\title{Homework 3, IEOR 6614}
	\author{Zongyi Liu}
	\date{Mon, Mar 2, 2026}
	\begin{document}
		\maketitle
		
		\section{Question 1}
		
		Suppose that at some point in the execution of a push-relabel algorithm, there exists an integer 
		$0 < k \le n - 1$ for which no vertex has $d(v) = k$. Show that all vertices with $d(v) > k$ are 
		on the source side of a minimum $s$--$t$ cut. If such a $k$ exists, the gap heuristic updates every 
		vertex $v$ for which $d(v) > k$ to set
		\[
		d(v) = \max(d(v), n + 1).
		\]
		Show that $d$ is still a valid distance function. (The gap heuristic is crucial in making 
		implementations of the push-relabel method perform well in practice.)
		
		\textbf{Answer}
		
		Let $A = \{v \in V : d(v) > k\}$.
		
		\textbf{Step 1: Vertices with $d(v) > k$ lie on the source side of a minimum cut.}
		
		First note that $d(s) = n > k$ and $d(t) = 0 \le k$. Hence
		$s \in A$ and $t \notin A$, so $A$ separates $s$ and $t$.
		
		Now consider any residual edge $(u,v)$ with $u \in A$. Since $d$ is a valid
		distance labeling, for every residual edge $(u,v)$ we have
		$d(u) \le d(v) + 1$. Thus
		$d(v) \ge d(u) - 1.$
		
		Since $u \in A$, we have $d(u) > k$, hence $d(u) \ge k+1$, and therefore
		$d(v) \ge k.$
		
		By assumption no vertex has label $k$. Hence $d(v) \ne k$ and therefore
		$d(v) > k$. Thus $v \in A$.
		
		Therefore no residual edge goes from $A$ to $V \setminus A$. Hence the cut
		$(A, V \setminus A)$ has residual capacity $0$. Since cut capacities are
		nonnegative, this cut is a minimum $s$--$t$ cut in the residual network.
		Thus every vertex with $d(v) > k$ lies on the source side of a minimum
		$s$--$t$ cut.
		
		\medskip
		
		\textbf{Step 2: The gap heuristic preserves validity of the distance labels.}
		
		The gap heuristic updates the labels as
		$
		d'(v) =
		\begin{cases}
			\max(d(v), n+1), & \text{if } d(v) > k, \\
			d(v), & \text{otherwise}.
		\end{cases}
		$
		
		We show that $d'$ remains a valid distance function. A valid labeling must
		satisfy for every residual edge $(u,v)$:
		$d'(u) \le d'(v) + 1.$
		
		We consider cases.
		
		\medskip
		
		\textit{Case 1: $u \notin A$.}
		
		Then $d'(u) = d(u)$ and $d'(v) \ge d(v)$. Since $d$ was valid,
		$d(u) \le d(v) + 1 \le d'(v) + 1.$
		Thus $d'(u) \le d'(v) + 1$.
		
		\medskip
		
		\textit{Case 2: $u \in A$.}
		
		From Step 1, if $(u,v)$ is residual and $u \in A$, then $v \in A$ as well.
		Thus both endpoints lie in $A$.
		
		Now $d'(u) = \max(d(u), n+1)$ and $d'(v) = \max(d(v), n+1)$.
		
		If $d'(u) = n+1$, then
		$d'(u) = n+1 \le d'(v) + 1.$
		
		If $d'(u) = d(u)$, then using validity of $d$,
		$d'(u) = d(u) \le d(v) + 1 \le d'(v) + 1.$
		
		Thus the inequality holds in all cases.
		
		\medskip
		
		Therefore $d'$ satisfies the distance labeling condition and remains a
		valid distance function.
		
		\medskip
		
		\textbf{Conclusion.}  
		If there exists a gap at label $k$, then all vertices with $d(v) > k$ lie
		on the source side of a minimum $s$--$t$ cut. Updating their labels to
		$\max(d(v), n+1)$ preserves the validity of the distance function, which
		justifies the correctness of the gap heuristic.
		
		\clearpage
	
	\section{Question 2}
	
	In class, we talked about efficient push-relabel algorithms using the following
	procedure for processing vertices. In this procedure $c_f$ is residual capacity
	and the neighbors of a vertex are stored in a linked list which has
	{current}, {head} and {next-neighbor} functions.
	
	\medskip

	
	\begin{algorithm}[h]
		\caption{\textsc{Discharge}$(u)$}
		\begin{algorithmic}[1]
			\While{$e(u) > 0$}
			\State $v \gets u.\mathit{current}$
			\If{$v = \mathrm{NIL}$}
			\State \textsc{Relabel}$(u)$
			\State $u.\mathit{current} \gets u.\mathit{head}$
			\ElsIf{$c_f(u,v) > 0$ \textbf{and} $d(u) = d(v) + 1$}
			\State \textsc{Push}$(u,v)$
			\Else
			\State $u.\mathit{current} \gets v.\mathit{next\mbox{-}neighbor}$
			\EndIf
			\EndWhile
		\end{algorithmic}
	\end{algorithm}
	
	\medskip
	
	Suppose that we run a push-relabel algorithm, but when it comes time to choose
	an operation, we always run discharge on the active vertex with the highest
	distance label. Show that this algorithm runs in $O(n^3)$ time.
	
	\medskip
	
	\emph{Hint:} Follow the analysis of the generic algorithm, but introduce a
	different potential function to help bound the number of non-saturating pushes.
	
		\textbf{Answer}
		
		We analyze the running time of the push--relabel algorithm that always
		selects the active vertex with the highest distance label and performs
		\textsc{Discharge} on it.
		
		\medskip
		
		First recall several known properties of the push--relabel algorithm.
		
		\begin{itemize}
			\item The label of every vertex is at most $2n-1$.
			\item Each relabel operation strictly increases $d(v)$.
			\item Therefore each vertex can be relabeled at most $2n$ times.
		\end{itemize}
		
		Since there are $n$ vertices, the total number of relabel operations is
		$O(n^2).$
		
		\medskip
		
		Next consider push operations. There are two types of pushes.
		
		\medskip
		
		\textit{Saturating pushes.}
		
		A push along $(u,v)$ is saturating if it makes the residual capacity
		$c_f(u,v)$ equal to zero. Each saturating push eliminates one residual
		edge. Since each edge can only become saturated a limited number of
		times between relabel operations, the total number of saturating pushes
		is
		$O(n^2).$
		
		\medskip
		
		\textit{Non-saturating pushes.}
		
		To bound the number of non-saturating pushes, define a potential function
		$
		\Phi = \sum_{v \text{ active}} d(v).
		$
		
		Since the algorithm always selects the active vertex with the highest
		label, a non-saturating push from $u$ to $v$ makes $u$ inactive and
		possibly activates $v$. Because admissible edges satisfy
		$d(u) = d(v) + 1,$
		the potential strictly decreases.
		
		Each relabel operation increases $\Phi$ by at most $n$, and since the
		total number of relabels is $O(n^2)$, the total increase of $\Phi$ is
		$O(n^3).$
		
		Because every non-saturating push decreases $\Phi$ by at least $1$, the
		total number of non-saturating pushes is also
		$O(n^3).$
		
		\medskip
		
		Finally, each execution of \textsc{Discharge} performs only constant work
		between pushes or relabels, so the total running time of the algorithm
		is bounded by the number of these operations. The total number of relabels is $O(n^2)$ and the total number of pushes
		is $O(n^3)$. Therefore the running time of the algorithm is
		$O(n^3).$

		
		\clearpage
		
		\section{Question 3}
		\begin{enumerate}
			\item[(a)] Give an algorithm with running time $O(n^2)$ which returns an order
			$v_1, v_2, \ldots, v_n$ of the vertices such that for all $i = 2,3,\ldots,n$,
			\[
			v_i
			=
			\arg\max_{w \in \{v_i,v_{i+1},\ldots,v_n\}}
			\sum_{\substack{vw \in E:\\ v \in \{v_1,\ldots,v_{i-1}\}}} u(vw).
			\]
			(More efficient algorithms that run in time $O(m + n\log n)$ are known for this problem).
			
			\item[(b)] Fix the order we obtained from the previous step. Prove that the minimum
			capacity of a cut separating $v_{n-1}$ and $v_n$ is
			\[
			\sum_{\substack{v v_n :\\ v \in \{v_1,\ldots,v_{n-1}\}}} u(vv_n).
			\]
			
			\item[(c)] Use the previous facts to deduce an algorithm which finds the minimum cut
			of $G$ in $O(n^3)$ time.
		\end{enumerate}
		
		
		\textbf{Answer}
		
		\item[(a)] We construct the order greedily. For each vertex $w$ not yet chosen, let
		$
		\alpha(w) = \sum_{\substack{vw \in E:\\ v \in \{v_1,\ldots,v_{i-1}\}}} u(vw),
		$
		that is, the total capacity from $w$ to the set of already chosen vertices.
		
		\medskip
		
		\textit{Algorithm.}
		Choose $v_1$ arbitrarily. For $i=2,3,\ldots,n$, having already chosen
		$v_1,\ldots,v_{i-1}$, select
		$
		v_i \in \arg\max_{w \in V \setminus \{v_1,\ldots,v_{i-1}\}} \alpha(w).
		$
		After selecting $v_i$, update the values $\alpha(w)$ for all unchosen neighbors
		$w$ of $v_i$ by adding $u(v_iw)$.
		
		\medskip
		
		A simple implementation stores the set of chosen vertices and, for each unchosen
		vertex, its current key $\alpha(w)$. At step $i$ we scan all unchosen vertices to
		find the maximum, which takes $O(n)$ time, and then update all keys in $O(n)$ time.
		Thus each step costs $O(n)$, and over $n$ steps the total running time is $O(n^2).$
		
		By construction, for every $i=2,\ldots,n$,
		$
		v_i
		=
		\arg\max_{w \in \{v_i,v_{i+1},\ldots,v_n\}}
		\sum_{\substack{vw \in E:\\ v \in \{v_1,\ldots,v_{i-1}\}}} u(vw).
		$
		
		\medskip
		
		\item[(b)] Let
		$
		A_i = \{v_1,\ldots,v_i\}.
		$
		We must show that the minimum capacity of a cut separating $v_{n-1}$ and $v_n$ is
		$
		\sum_{\substack{v v_n \in E:\\ v \in \{v_1,\ldots,v_{n-1}\}}} u(vv_n).
		$
		
		Let $(S,\overline S)$ be any cut with $v_n \in S$ and $v_{n-1} \in \overline S$.
		Let $j$ be the smallest index such that $v_j \in S$. Then
		$
		v_1,\ldots,v_{j-1} \in \overline S.
		$
		Since $v_j$ was chosen by maximum adjacency, we have
		$
		\sum_{\substack{v v_j \in E:\\ v \in A_{j-1}}} u(vv_j)
		\ge
		\sum_{\substack{v v_n \in E:\\ v \in A_{j-1}}} u(vv_n).
		$
		Now the left-hand side counts only edges crossing the cut from $v_j$ to
		$\overline S$, so it is at most the capacity of the cut $(S,\overline S)$.
		Hence
		$
		u(S,\overline S)
		\ge
		\sum_{\substack{v v_n \in E:\\ v \in A_{j-1}}} u(vv_n).
		$
		Since $A_{j-1} \subseteq \{v_1,\ldots,v_{n-1}\}$ and all capacities are nonnegative,
		we get
		$
		u(S,\overline S)
		\ge
		\sum_{\substack{v v_n \in E:\\ v \in \{v_1,\ldots,v_{n-1}\}}} u(vv_n).
		$
		
		On the other hand, the cut
		$
		(\{v_n\}, V \setminus \{v_n\})
		$
		separates $v_{n-1}$ and $v_n$, and its capacity is exactly
		$
		\sum_{\substack{v v_n \in E:\\ v \in \{v_1,\ldots,v_{n-1}\}}} u(vv_n).
		$
		Therefore this is the minimum capacity of any cut separating $v_{n-1}$ and $v_n$.
		
		\medskip
		
		\item[(c)] The previous parts give the standard contraction algorithm for global minimum cut.
		
		\medskip
		
		\textit{Algorithm.}
		Repeatedly do the following until only one vertex remains:
		
		\begin{enumerate}
			\item Compute a maximum-adjacency ordering
			$
			v_1,\ldots,v_n
			$
			in the current graph, using part (a).
			\item Let
			$
			\lambda = \sum_{\substack{v v_n \in E:\\ v \in \{v_1,\ldots,v_{n-1}\}}} u(vv_n).
			$
			By part (b), this is the minimum capacity of a cut separating $v_{n-1}$ and $v_n$.
			Record this value.
			\item Contract the two vertices $v_{n-1}$ and $v_n$ into a single vertex, combining
			parallel edges by adding their capacities.
		\end{enumerate}
		
		Return the smallest recorded value over all phases.
		
		\medskip
		
		\textit{Correctness.}
		At each phase, part (b) gives an $s$--$t$ minimum cut for the pair
		$(v_{n-1},v_n)$. If a global minimum cut separates these two vertices, then we
		have found its value. If not, then $v_{n-1}$ and $v_n$ lie on the same side of
		some global minimum cut, so contracting them preserves that cut. Therefore some
		phase must record the value of a global minimum cut, and taking the minimum over
		all phases returns the global minimum cut value.
		
		\medskip
		
		\textit{Running time.}
		Each phase computes the ordering in $O(n^2)$ time and then performs one contraction.
		There are $n-1$ phases, so the total running time is
		$
		O(n^3).
		$
		
		
		\clearpage
		
		\section{Question 4}
		
		\begin{enumerate}
			\item[(a)] Use the analysis of the Contraction Algorithm for global minimum cuts
			to prove that a graph with $n$ vertices can have at most $\binom{n}{2}$
			distinct minimum cuts (distinct up to complement).
			
			\item[(b)] Provide an example of a specific graph topology on $n$ vertices that has
			exactly $\binom{n}{2}$ minimum cuts, proving that the analysis of Karger's
			algorithm is tight.
		\end{enumerate}
	
	
	\textbf{Answer}
	
	\begin{enumerate}
		\item[(a)] Let $\lambda$ denote the value of a global minimum cut of $G$.
		We use the standard analysis of Karger's contraction algorithm.
		
		Fix any particular minimum cut $C$. In the contraction algorithm, as long as
		no edge of $C$ is contracted, the two sides of $C$ remain distinct, and so
		$C$ survives to the end. Thus the algorithm outputs $C$ provided it never
		contracts an edge of $C$.
		
		Now suppose that at some stage of the algorithm there are $r$ vertices left.
		Because the minimum cut value is $\lambda$, every vertex has degree at least
		$\lambda$, so the graph has at least
		$
		\frac{r\lambda}{2}
		$
		edges. Since the cut $C$ has exactly $\lambda$ edges, the probability that the
		next random contraction chooses an edge of $C$ is at most
		$
		\frac{\lambda}{r\lambda/2} = \frac{2}{r}.
		$
		Hence the probability that $C$ survives this contraction is at least
		$
		1-\frac{2}{r}=\frac{r-2}{r}.
		$
		
		Therefore the probability that $C$ survives all contractions from $n$ vertices
		down to $2$ vertices is at least
		$
		\prod_{r= n}^{3}\frac{r-2}{r}
		=
		\frac{n-2}{n}\cdot\frac{n-3}{n-1}\cdots\frac{1}{3}
		=
		\frac{2}{n(n-1)}
		=
		\frac{1}{\binom{n}{2}}.
		$
		
		So every fixed minimum cut is output by the contraction algorithm with
		probability at least $1/\binom{n}{2}$.
		
		Now let $N$ be the number of distinct minimum cuts, counted up to complement.
		The algorithm outputs exactly one cut, so the probabilities of outputting
		these distinct minimum cuts sum to at most $1$. Since each one is output with
		probability at least $1/\binom{n}{2}$, we obtain
		$
		N \cdot \frac{1}{\binom{n}{2}} \le 1.
		$
		Thus
		$
		N \le \binom{n}{2}.
		$
		
		Hence a graph on $n$ vertices has at most $\binom{n}{2}$ distinct minimum cuts
		(up to complement).
		
		\medskip
		
		\item[(b)] Consider the cycle graph $C_n$ on vertices
		$
		v_1,v_2,\ldots,v_n,
		$
		with unit capacity on every edge.
		
		The value of the global minimum cut is $2$, since every vertex has degree $2$,
		and isolating any nonempty proper set of consecutive vertices cuts exactly two
		edges.
		
		Now every minimum cut in $C_n$ is obtained by choosing two distinct edges of
		the cycle and deleting them. Removing two distinct edges breaks the cycle into
		two connected components, which define a cut of capacity $2$. Conversely, any
		cut of capacity $2$ must correspond to removing exactly two edges, since all
		edge capacities are $1$ and no cut can have capacity $1$ in a cycle.
		
		Thus the minimum cuts are in one-to-one correspondence with unordered pairs of
		distinct edges of the cycle. Since the cycle has $n$ edges, the number of such
		pairs is
		$
		\binom{n}{2}.
		$
		
		Two complementary sides define the same cut, and choosing the same two deleted
		edges gives exactly one cut up to complement, so the number of distinct minimum
		cuts (up to complement) is exactly
		$
		\binom{n}{2}.
		$
		
		Therefore the bound in part (a) is tight.
	\end{enumerate}
	
		\clearpage
		
		\section{Question 5}
		Consider the following modification of the maximum flow problem. You are given:
		a network $N(V,A)$; $k$ goods and for each good $i$ a source $s_i$ and a sink
		$t_i$; and upper bounds $u_a$ on each arc $a$. A feasible flow is a collection
		of non-negative vectors $f_1,\ldots,f_k \in \mathbb{R}^A$ such that each $f_i$
		is an $s_i$--$t_i$ flow, and such that, on each arc $a$, the total flow passing
		through that arc is at most $u_a$. The problem is to find an integer flow that
		maximizes $\sum_{i=1}^k \operatorname{val}(f_i)$, where $\operatorname{val}(f_i)$
		is the value of a flow, that is the flow leaving the source (or equivalently,
		entering the sink).
		
		\begin{enumerate}
			\item[(a)] Write the problem as an Integer Program.
			
			\item[(b)] Suppose that all data is integer, nonnegative, and that the problem is
			feasible. Relax the integrality constraints of your IP, and consider the LP you
			obtain. Does this problem always have an optimal solution that is integer and why?
		\end{enumerate}
	
		\textbf{Answer}
		
		\textbf{Answer.}
		
		\begin{enumerate}
			\item[(a)] We introduce variables $f_i(a)$ for each good $i=1,\ldots,k$ and each
			arc $a \in A$, where $f_i(a)$ denotes the amount of good $i$ sent through arc $a$.
			We also introduce a variable $x_i$ for the value of the $i$th flow.
			
			The integer program is:
			
			$
			\max \sum_{i=1}^k x_i
			$
			
			subject to the following constraints.
			
			\medskip
			
			\textit{Flow conservation for each commodity.} For every good $i$ and every
			vertex $v \in V \setminus \{s_i,t_i\}$,
			$
			\sum_{a \in \delta^+(v)} f_i(a) - \sum_{a \in \delta^-(v)} f_i(a) = 0.
			$
			
			At the source $s_i$,
			$
			\sum_{a \in \delta^+(s_i)} f_i(a) - \sum_{a \in \delta^-(s_i)} f_i(a) = x_i.
			$
			
			At the sink $t_i$,
			$
			\sum_{a \in \delta^+(t_i)} f_i(a) - \sum_{a \in \delta^-(t_i)} f_i(a) = -x_i.
			$
			
			\medskip
			
			\textit{Arc-capacity constraints.} For every arc $a \in A$,
			$
			\sum_{i=1}^k f_i(a) \le u_a.
			$
			
			\medskip
			
			\textit{Nonnegativity and integrality.} For every $i$ and every $a \in A$,
			$
			f_i(a) \ge 0, \qquad f_i(a) \in \mathbb{Z},
			$
			and
			$
			x_i \ge 0, \qquad x_i \in \mathbb{Z}.
			$
			
			This is an integer program for the stated problem.
			
			\medskip
			
			\item[(b)] Relaxing the integrality constraints gives a linear program. In general,
			this LP does \emph{not} always have an optimal solution that is integer.
			
			The reason is that this is a \emph{multicommodity flow} problem. For ordinary
			single-commodity maximum flow, the constraint matrix is totally unimodular, so
			with integer capacities every extreme point solution is integral. Here, however,
			the different commodities are coupled through the shared arc-capacity constraints
			$
			\sum_{i=1}^k f_i(a) \le u_a,
			$
			and the resulting constraint matrix is not, in general, totally unimodular.
			Therefore fractional optimal solutions can occur.
			
			\medskip
			
			A simple example is the following. Let there be two goods, each of which wants to
			send one unit across a network in which the two commodities must share a bottleneck
			of capacity $1$. Fractionally, the LP may route $\tfrac12$ unit of each commodity
			through the bottleneck and achieve total value $1$, while no integral solution can
			do this for both simultaneously in the same way. Thus the LP relaxation need not
			have an integral optimum.
			
			Hence the answer is:
			
			\medskip
			
			\textbf{No,} the LP relaxation does not always have an optimal solution that is
			integer, because multicommodity flow does not in general enjoy the integrality
			property of single-commodity flow.
		\end{enumerate}
	
		
		
	\end{document}
